\subsection{Golden Point Nexic Sorting}

\begin{prince}
    Consider Nexus $A_1, A_2, \cdots A_n$ which serve an Aethic partition of Nexus $A = \bigcup_i A_i$, and consider the Golden Aethus $G$ and Cosmic Nexus $N_0$. The probability of an Aetherian inhabiting Aethus $G \cap A_i$ is proportional to the probability of $G$ and $A_i$ relative to $A$ via the Cosmic Nexus. 
    \begin{equation}
        \boxed{
        \operatorname{P} \! \left( G \cap A_i \;\middle|\; A, H_{\textnormal{in}} \right) \propto \operatorname{P} \! \left( G \cap A_i \;\middle|\; A, N_0 \right)
        }
    \end{equation}
    In other words, the Aethic weighting scheme for a given Aetherian coming to inhabit a given Aethus $A_i$ is such that this inhabitance and the Golden Point both occurring is only ever as likely as the probability of the Golden Point and $A_i$ occurring from $A$ via the Cosmic Nexus. 

    Very simply put, then, this is a fundamental assessment that the very probability of an Aetherian coming to inhabit a given Aethus is directly and intimately tied to the probability of the Golden Point occurring under the presence of that Aethus as opposed to a given alternative. 
\end{prince}

By the law of conditional probability, let us consider the following identity. Note that this assumes that once an Aetherian has come to inhabit Nexus $A_i$, all further centric unfolding will abide by the Aethic weighting scheme of $A_i$ itself. 
\begin{equation}
    \operatorname{P} \! \left( G \cap A_i \;\middle|\; A, H_{\textnormal{in}} \right) = \operatorname{P} \! \left( G \;\middle|\; A_i \right) \operatorname{P} \! \left( A_i \;\middle|\; A, H_{\textnormal{in}} \right)
\end{equation}
According to this, then, we have that the total probability of the Golden Point, by every avenue of $A_i$, may be expressed as a sum of all such contributions for every $A_i$, (that is, since $A_i$ are an Aethic partition of $A$, and so are disjoint).
\begin{equation}
    \begin{aligned}
        \operatorname{P} \! \left( G \;\middle|\; A, H_{\textnormal{in}} \right) & = \sum_i \operatorname{P} \! \left( G \cap A_i \;\middle|\; A, H_{\textnormal{in}} \right) \\
        & = \sum_i \operatorname{P} \! \left( G \;\middle|\; A_i \right) \operatorname{P} \! \left( A_i \;\middle|\; A, H_{\textnormal{in}} \right)
    \end{aligned}
\end{equation}
Now let us consider a simple example where $A$ is simply Aethically partitioned into $A_1$ and $A_2$, such that $A_1$ occurs directly prior to Miracle $M$, and $A_2$ occurs directly subsequently to it, (at least in the sense of Aethic time, where the knowledge or ability for deduction of its success is unknown versus known, respectively). 
\begin{equation}
    \dfrac{\operatorname{P} \! \left( G \cap A_1 \;\middle|\; A, H_{\textnormal{in}} \right)}{\operatorname{P} \! \left( G \cap A_2 \;\middle|\; A, H_{\textnormal{in}} \right)} = \dfrac{\operatorname{P} \! \left( G \;\middle|\; M \right) \operatorname{P} \! \left( M \;\middle|\; A_1 \right) \operatorname{P} \! \left( A_1 \;\middle|\; A, H_{\textnormal{in}} \right)}{\operatorname{P} \! \left( G \;\middle|\; A_2 \right) \operatorname{P} \! \left( A_2 \;\middle|\; A, H_{\textnormal{in}} \right)}
\end{equation}
With $A_2$ being implied to effectively occur directly after the Miracle in Aethic time, we have that $\operatorname{P} \! \left( G \;\middle|\; M \right) \approx \operatorname{P} \! \left( G \;\middle|\; A_2 \right)$, such that the further treks still needed through to the Golden Point is equivalent, and thereby probabalistically equal. This allows us to simplify further. 
\begin{equation}
    \dfrac{\operatorname{P} \! \left( G \cap A_1 \;\middle|\; A, H_{\textnormal{in}} \right)}{\operatorname{P} \! \left( G \cap A_2 \;\middle|\; A, H_{\textnormal{in}} \right)} = \dfrac{\operatorname{P} \! \left( A_1 \;\middle|\; A, H_{\textnormal{in}} \right)}{\operatorname{P} \! \left( A_2 \;\middle|\; A, H_{\textnormal{in}} \right)} \operatorname{P} \! \left( M \;\middle|\; A_1 \right)
\end{equation}
Likewise, in the case of the Cosmic Nexus-involving probabilities, we have that 
We may accordingly write the following using the statement of the principle itself, specifically through use of equality between the given partial probabilities. 
\begin{equation}
    \begin{aligned}
        \dfrac{\operatorname{P} \! \left( G \cap A_1 \;\middle|\; A, H_{\textnormal{in}} \right)}{\operatorname{P} \! \left( G \cap A_1 \;\middle|\; A, N_0 \right)} & = \dfrac{\operatorname{P} \! \left( G \cap A_2 \;\middle|\; A, H_{\textnormal{in}} \right)}{\operatorname{P} \! \left( G \cap A_2 \;\middle|\; A, N_0 \right)}, \\
        \dfrac{\operatorname{P} \! \left( G \cap A_1 \;\middle|\; A, H_{\textnormal{in}} \right)}{\operatorname{P} \! \left( G \cap A_1 \;\middle|\; A, H_{\textnormal{in}} \right) + \operatorname{P} \! \left( G \cap A_2 \;\middle|\; A, H_{\textnormal{in}} \right)} & = \dfrac{\operatorname{P} \! \left( G \cap A_1 \;\middle|\; A, N_0 \right)}{\operatorname{P} \! \left( G \cap A_1 \;\middle|\; A, N_0 \right) + \operatorname{P} \! \left( G \cap A_2 \;\middle|\; A, N_0 \right)} \\
        \dfrac{1}{1 + \dfrac{\operatorname{P} \! \left( G \cap A_2 \;\middle|\; A, H_{\textnormal{in}} \right)}{\operatorname{P} \! \left( G \cap A_1 \;\middle|\; A, H_{\textnormal{in}} \right)}} & = \dfrac{1}{1 + \dfrac{\operatorname{P} \! \left( G \cap A_2 \;\middle|\; A, N_0 \right)}{\operatorname{P} \! \left( G \cap A_1 \;\middle|\; A, N_0 \right)}} \\
        
    \end{aligned}
\end{equation}
By this example principle, then, we may write 

\section{Retrocausality}

\begin{prince}[Retrocausal Accordance Principle]
    For a given Aetherian, whose Aethus is $H_1$, and a given stated-attribute in their Aethus, $A \supset H_1$ such that $A$ is extremely unlikely in the Contentor, $N_0$, we have that the probability of said Aetherian centric unfolding to a future Aethus containing the stated-attribute $B$ is bounded from above by the probability of $A$ with respect to $B$ and the Contentor. 
    \begin{equation}
        \left. \mathrm{P} \! \left( B \mid H_1, N_0 \right) \lesssim \mathrm{P} \! \left( A \mid B, N_0 \right) \, \middle| \, H_1 \subset A, \mathrm{P} \! \left( A \mid N_0 \right) \ll 1 \right.
    \end{equation}
    Note that this is presumed true for an arbitrary Contentor $N_0$. 

    Derivationally speaking, this principle is a direct corollary of the anthropic accordance principle \cite{Nexus}, specifically where we consider how the state of $A$ being exceptionally unlikely with respect to one's future Aethus $H_1 \cap B$ then implies that one must have witnessed a fluke with respect to that Aethus by even currently being present in the Aethus $H_1$ in the first place, (for which $A$ has occurred). Accordingly, then, we have that this consistences a fluke, (to the extent of the anthropic accordance principle's accuracy rate), if said current Aethus dissatisfies this mediocrity constraint with respect to any one of one's future Aethae, (since after all said Aethae will have to eventually be relevant to oneself, so as to then warrant the application of the first postulate of the Aethus to them on one's behalf). So, we conclude that the probability of so much as centric unfolding to these dubious futures in the first place must be constrained at least proportionally to their degree of dubiousness, (so to speak). 

    In short, the simplified motto of this principle might be that ``your present cannot be arbitrary to your future". 
\end{prince}

Allow me to provide a direct example of my applying this principle to my day-to-day life, which might hopefully help with the general analogy. So, when I was in grade school, I was told on a few occasions that I have an exceptionally strong presentational speaking ability. Having this in my mind while I was deriving Nexic reasoning in 2021, I then had a rather intriguing thought: \textit{if I go my entire life without ever needing the skill of speaking ability, then would this not imply that my state of having this skill right now in the first place represents an immense fluke with respect to said life?} That is, if I imagine the set of all lives I could live in which I never need to be good at presenting, then living this current life of mine where I am good at presenting would be a decided fluke in the set such lives. So, therefore, as an effective proof by contradiction needed to disallow this waking moment from being a fluke, I would need to infer that somewhere, somehow in my future I will have to complete at least one task that requires a strong speaking ability, (thereby satisfying the motto of the retrocausal accordance principle by keeping my present very much non-arbitrary with respect to my future). So, somehow or another, I will have constrained a future task of mine just under the assumption that I am not currently living an anomaly in the Aethus of all possible lives I could have led. Perhaps the most remarkable effect of this idea is that I seem to have been right, simply because I now know that I needed a strong speaking ability in order to write and present these very papers of mine which you now read. In effect, I will have predicted an attribute of a future Aethus of mine by simply analyzing my own standing fluke content. Altogether, the broad conclusion becomes the form of the retrocausal accordance principle itself, for which all Aetherians may similarly constrain the probabilities of their future events given their current Aethus of fluke events. So, to be abundantly clear, the idea is a variation of the Accordance Principle where instead of conditioning on past light cone events so that the current moment stays out of being a fluke, we instead condition on whichever future Aethic time Aethae make the present moment minimally a fluke, (but with the graded $\lessim$ so as to allow edge cases, of course).

Also for another major part: by the concept of Golden Optimation, you personally will have just such fluke traits if the Golden Point conditioning itself required you having them to significantly help its existence. So in this sense the Golden Point conditioning would be the real first-principles deduction explaining \textit{why} you have those traits, whereas the retrocausal accordance principle being applied to that the traits will probably \textit{pay off} in the future is just inference rather than first principles-deduction. Also note that the converse does not generally hold to this, meaning that just because you have a fluke does not mean Golden Optimation triggered it, but it does hold that if the Golden Point required you having the fluke trait, then conditioning on the Golden Point implies that you must have it. In the Golden Route World itself, we then attain a profound shift in how causality and happenstance themselves work, basically as a kind of ultimate effect of Aethic reasoning rendering causality as emergent and then being able to introduce catalysts to the Aethic structure which enact wild swings to the behavior of causality: for a given notoriously challenging feat, under strong extrapolation, the leap of faith is not “\textit{will I achieve this?}” but “\textit{is this genuinely required for the Golden Point, and that I'm the one to do it?}”. It becomes a different question because of Golden Optimation, but it remains a leap of faith because of the risk that it might not be necessary for the Golden Point at all, in which case even in the Route World it is not probabilistically privileged as an Aethic event. 

\subsection{Retrocausal Argument for the Golden Point}

Perhaps the central argument for why the Golden Point ought to be taken seriously as the solution to the finale paradox comes from the retrocausal accordance principle itself. 
\begin{concept}[Retrocausal Argument for the Golden Point]
    Let us for a moment consider the totality of all Aethic stated-attributes about our biosphere which, relative to one's Aethus, apply exclusively to one's living in the present day as opposed to during any other stretch of time throughout Earth's natural history. There might be the fact that the hottest \cite{ALICE2012Temperature} and coldest \cite{Kovachy2015Picokelvin} temperatures in the observable universe exist on our own planet, that the first vertebrate to pass the Kármán line was born within the last hundred years \cite{burgess2007animals}, or that unlocking dormant genes by genetic engineering is suddenly possible \cite{Pattali2026_CRISPR_epigenome_mod}. Let us suppose that we refer to the Aethus $A$ as the Aethus of all of these events occurring, with arbitrary parent Aethus $A_i$ to this Aethus essentially representing a respectively arbitrary ``subset" of them occurring. 
    By the retrocausal accordance principle, now, we might infer for each such $A_i$ that there is more likely than not to exist some stated-attribute $B_i$ in the future of our biosphere for which $A_i$ is almost surely a necessary requisite as opposed to the alternative of one's having been born far further back into deep time. What we get now is essentially something of an optimization query, being that for an $A_i$ of only a fraction of the total developments, for there not to exist future stated-attribute $B_i$ to not take any other of the developments as necessary is quite unlikely on account of such likely violating the retrocausal accordance principle, yet on the other hand for the entirety of developments in $A$ there ought to perhaps be at least a handful of developments which are not necessitated by future stated-attributes, since there is probabilistic margin of error to the retrospective accordance principle, however slight. Altogether then the estimate is that there ought to exist some $A_j$ for which the significant majority of the developments in $A$ all root to a future stated-attribute which may captured concisely in the single Aethus $B_j$. Intuitively, this ought to have the visualization of a most extremely connected graph, where this one single optimized $B_j$ serves to essentially act as a single future stated-attribute which takes each and every modern development in $A_j$ as a total necessity for its occurrence. The philosophical question, then, becomes this: \emph{contextually speaking, what might we infer such a $B_j$ to look like?} This is the matter for which we may establish two major points of argument. 
    \begin{enumerate}
        \item Such a constrained future stated-attribute $B_j$, remember, ought to render all modern developments absolutely probabilistically and categorically necessary as prerequisites against the backdrop of their not having yet occurred across the expanses of deep time. So what we are looking for in predicting the shape of $B_j$ which not only correlates strongly with the totality of all human progress in manipulating our environment, but is essentially exactly the stated-attribute with the strongest possible correlation to such things.
        \item This optimized correlation is because the full probabilistic degree to which $B_j$ fails to be correlated with $A_j$ is essentially an opening through which one's holding a farther back Aethus in deep time would be compatible with $B_j$. So if $B_j$ itself ought to directly source the reasoning for why one finds oneself in the modern day rather than natural history via the retrocausal accordance principle, then we might make the following probabilistic argument with it as the reference. Suppose the Aetherian Nexus is defined as $N_0 \cap B_j$, where $N_0$ is the Cosmic Nexus. We then gather that the probability of oneself existing in the modern day as opposed to farther back in deep time is the following.
        \begin{equation}
            P \! \left( A_j \mid N_0 \cap B_j \right) = \frac{P \! \left( A_j \cap B_j \mid N_0 \right)}{P \! \left( A_j \cap B_j \mid N_0 \right) + P \! \left( A_j' \cap B_j \mid N_0 \right)}
        \end{equation}
        If we now consider essentially Aethic partitioning $B_j$ into $C$, a component of extreme correlation to $A_j$, and $C'$, which is roughly statistically independent to $A_j$, when we might have some linear combination of $B_j = \left( 1 - \alpha \right) C + \alpha C'$. Accordingly, then, by intersecting $A_j' \cap B_j$ we have that the first term falls out due to almost sure invalidity, and so the expression becomes $A_j' B_j = \alpha A_j' C'$. Now, because of the approximate statistical independence between $A_j'$ and $C'$, we may write the total probability of $A_j' \cap B_j \mid N_0$ as a product of probabilities. 
        \begin{equation}
            P \! \left( A_j \mid N_0 \cap B_j \right) = \frac{P \! \left( A_j \cap B_j \mid N_0 \right)}{P \! \left( A_j \cap B_j \mid N_0 \right) + \alpha P \! \left( A_j' \mid N_0 \right) P \! \left( C' \mid N_0 \right)}
        \end{equation}
        Now, we might consider the true vastness of the relative magnitude of $P \! \left( A_j' \mid N_0 \right)$ compared to $P \! \left( A_j \mid N_0 \right)$ through use of a simple exercise in exponentials. That is, if we were picking a random time in Earth's natural history over a uniform probability density function alone, then there would be about a ${10}^{-8}$ probability of landing in the last century, (being that which is necessary for $A_j$ to occur as opposed to the alternative), but this does not even take into account the exponential compounding of the internal and external doomsday dangers \cite{Nexus}. In effect, then, for net doomsday danger hazard rate $\lambda$, we have that already after time $\frac{\ln \! \left( 2 \right)}{\lambda}$ it is more likely than not that the landing would have occurred, because the Aethic weight of survival for the biosphere drops rapidly to near-zero afterwards. With $P \! \left( C' \mid N_0 \right) \approx 1$ perhaps being true within an order of magnitude or two due to its statistical independence to $A_j$, (as anyway most likely produces a random degree of probabilistic favorability within the Cosmic Nexus), what we are essentially left with is the following relationship, given that $P \! \left( A_j \mid N_0 \cap B_j \right) \gtrsim P \! \left( A_j' \mid N_0 \cap B_j \right)$ is essentially demanded by the retrocausal accordance principle.
        \begin{equation}
            \alpha \lesssim \frac{P \! \left( A_j \mid N_0 \right)}{P \! \left( A_j' \mid N_0 \right)} \approx 0
        \end{equation}
        Essentially, then, this is a rough explanation of why the degree to which $B_j$ ought to optimally require $A_j$ with a ferocity that is almost incomprehensible---even the slightest perturbations away from this optimal necessity, as would for a given choice of conventionalized difference inadvertently heighten the value of $\alpha$ to far beyond its optimally-selected location, would for all intents and purposes lead $P \! \left( A_j \mid N_0 \cap B_j \right)$ to approach zero with the full cosmic rapidity implied by the magnitude of $\frac{P \! \left( A_j' \mid N_0 \right)}{P \! \left( A_j \mid N_0 \right)}$. 
        
        In essence, then, we are not only looking for a $B_j$ in the future which satisfies the rendering of $A_j$ as necessary compared with all times prior, but one which on sight is unmistakably optimized at a cosmic level of significance.
    \end{enumerate}
    The argument for the Golden Point itself, then, is that just such an Aethus---with its particular implications for the indefinite survivorship of our biosphere---satisfies perhaps more strongly than any other contender the possibility of being just such a cosmically-significant $B_j$ Aethus. In a sense, it is exactly what we would expect such a causal entity to look like, because any other alternative does not quite present as ``ambitious" enough to satisfy this unmistakability condition. 

    So if you will, this argument not only solves the finale paradox, but uses its own momentum to argue for the significance of the Golden Point on the Aethic Tree.
\end{concept}
